% Paper 5: Carbon-Verify — LLM Testbench Generation for Hardware Verification
% Target: DVCon / DAC Workshop

\documentclass[conference]{IEEEtran}
\usepackage{amsmath,graphicx,booktabs,hyperref}

\title{Can LLMs Verify Hardware?\\Testbench Generation from Code vs.\ Specification}

\author{
\IEEEauthorblockN{William Chen}
\IEEEauthorblockA{National Taiwan University}
}

\begin{document}
\maketitle

\begin{abstract}
We investigate whether local small language models (4.5B parameters)
can generate functional hardware testbenches. Using Gemma4 with
iverilog simulation, we test two approaches: \emph{code-driven}
(generate testbench from module implementation) and \emph{spec-driven}
(generate from behavioral specification). Code-driven testbenches
achieve 100\% compile rate (3/3 modules) but catch 0\% of injected
bugs---because the model derives expected values from the buggy code
rather than the specification. Spec-driven testbenches fail to compile
(0/3), as the model struggles to translate natural language
specifications into syntactically valid Verilog. We identify a
fundamental limitation: \emph{LLM testbenches verify what code does,
not what code should do.} Cocotb (Python-based) testbenches show
higher success rates than pure Verilog, suggesting Python is a more
natural target language for LLM-generated verification.
\end{abstract}

\section{Introduction}

Hardware verification consumes 60--70\% of semiconductor design time.
Testbench generation---writing stimulus and checking logic---is
labor-intensive yet follows patterns that language models might learn.

We examine whether Gemma4 (4.5B MoE) can generate testbenches that
actually compile, simulate, and detect bugs in Verilog modules.
Two approaches are compared:

\begin{itemize}
    \item \textbf{Code-driven}: The model sees the module's RTL code
    and generates a testbench from it.
    \item \textbf{Spec-driven}: The model receives only a natural
    language specification of expected behavior, plus port declarations.
\end{itemize}

\section{Experimental Setup}

\subsection{Modules Under Test}
Three modules with intentionally injected bugs:
\begin{enumerate}
    \item \textbf{Counter}: Reset sets count to 1 instead of 0
    \item \textbf{Adder}: Adds extra +1 to sum
    \item \textbf{Mux4}: No bug (control case)
\end{enumerate}

\subsection{Tools}
iverilog for compilation, vvp for simulation, Gemma4 via Ollama.

\section{Results}

\begin{table}[t]
\centering
\caption{Testbench generation results. Code-driven compiles but
misses bugs. Spec-driven fails to compile.}
\label{tab:verify}
\begin{tabular}{lcccc}
\toprule
& \multicolumn{2}{c}{\textbf{Code-Driven}} & \multicolumn{2}{c}{\textbf{Spec-Driven}} \\
\cmidrule(lr){2-3} \cmidrule(lr){4-5}
\textbf{Module} & Compile & Bug? & Compile & Bug? \\
\midrule
Counter (bug)    & \checkmark & $\times$ & $\times$ & --- \\
Adder (bug)      & \checkmark & $\times$ & $\times$ & --- \\
Mux4 (correct)   & \checkmark & PASS     & ---      & --- \\
Shift Reg (bug)  & ---        & ---      & $\times$ & --- \\
\midrule
\textbf{Rate}    & 3/3 (100\%) & 0/2 (0\%) & 0/3 (0\%) & --- \\
\bottomrule
\end{tabular}
\end{table}

\subsection{Why Code-Driven Misses Bugs}

The counter testbench expected \texttt{count=1} after reset---matching
the buggy code, not the specification. The model learned expected
values by \emph{simulating} the code in its ``head,'' not by
reasoning about intended behavior.

This reveals a fundamental limitation:
\begin{quote}
\emph{LLM-generated testbenches from code verify that code does
what it does, not what it should do.}
\end{quote}

\subsection{Why Spec-Driven Fails to Compile}

Spec-driven generation requires translating natural language
(``reset should set count to 0'') into valid Verilog syntax.
Gemma4 struggles with:
(a) Verilog timing constructs (\texttt{\#}, \texttt{@}),
(b) correct port binding, and
(c) \texttt{\$display}/\texttt{\$finish} placement.

\subsection{Cocotb vs. Verilog}

Cocotb (Python-based) testbenches from our PoC experiments achieve
3/3 generation success (vs. 0/3 for spec-driven Verilog). Python's
simpler syntax and the model's stronger Python training data
make it a more reliable target language.

\section{Discussion}

\subsection{The Spec-Code Gap}

To generate bug-catching testbenches, the model needs \emph{both}:
(a) the specification (what should happen), and (b) enough Verilog
skill to express assertions. Current 4.5B models have (b) weakly
but lack the ability to bridge from (a) to (b).

\subsection{Toward Effective LLM Verification}

\begin{enumerate}
    \item \textbf{Use Python (Cocotb)} instead of Verilog for
    testbench generation
    \item \textbf{Provide explicit assertions} in the prompt
    (``assert count==0 after reset'')
    \item \textbf{Two-stage}: LLM generates test vectors,
    separate tool checks assertions
\end{enumerate}

\section{Conclusion}

LLMs can generate compilable testbenches (100\% code-driven) but
cannot currently detect bugs because they derive expectations from
code rather than specifications. Spec-driven generation is the
correct approach but fails on Verilog syntax (0\% compile rate).
Python-based frameworks (Cocotb) are the pragmatic path forward
for LLM-assisted verification.

\bibliographystyle{IEEEtran}
\begin{thebibliography}{1}
\bibitem{cocotb} Cocotb Contributors, ``Cocotb: Coroutine-based Co-simulation Testbench,'' 2024.
\end{thebibliography}

\end{document}
